\chapter{MOLAP Cubes}
\label{ch:molap}

Business rules often need aggregates: the total exposure of an account,
the average response time of an application, the number of orders per
region. Computing them in rules is awkward, because a rule sees one
combination of facts at a time. \Morendo\ instead lets you declare a
multidimensional cube that the engine keeps up to date as facts come and
go, and lets rules query the cube as if it were a fact. The design is
described in \file{doc/molap\_for\_production\_system.pdf}; this chapter
covers how to use it.

OLAP (online analytical processing) tools organize data the same way for
reporting: \emph{dimensions} are the attributes the data is grouped by,
such as region or product, and \emph{measures} are the numbers computed
for each group, such as a sum or an average. A conventional OLAP server
answers queries on demand, so an application that must react to a
changed figure has to poll it. A \Morendo\ cube is \emph{reactive}
instead: it lives in memory beside working memory, the engine updates it
as its source facts change, and the rules that read it are re-evaluated
when it changes, so the rules themselves decide which figures matter.

\section{Declaring a cube}

A cube has three parts, in this order: the \emph{sources}, one or more
patterns joined like the left-hand side of a rule; the \emph{dimensions},
the slots the cube can be sliced by; and the \emph{measures}, aggregates
computed over each slice. \file{samples/molap/cubes.clp} declares a cube
over web-server and database logs:

\begin{lstlisting}
(defcube TransactionCube
  <- (HTTPLog
    (IPAddress ?ipaddr)
    (URL ?url)
    (bytesSent ?bytesSent)
    (elapsedTime ?elapsedTime)
    (endTime ?endTime)
    (serverType ?applicationType)
    (statusCode ?statusCode)
    (timestamp ?timestamp)
    (transactionID ?transactionId)
  )
  <- (DatabaseQueryLog
    (IPAddress ?dbipaddr)
    (bytesSent ?dbBytesSent)
    (elapsedTime ?dbElapsedTime)
    (endTime ?dbEndTime&:(within-minutes 4 ?dbEndTime ?endTime) )
    (statusCode ?dbStatusCode)
    (timestamp ?dbTimestamp&:(within-minutes 2 ?timestamp ?dbTimestamp) )
    (transactionID ?transactionId )
  )
  (dimension IPAddress ?ipaddr)
  (dimension Application ?url)
  (dimension TransactionID ?transactionId)
  (dimension ServerType ?applicationType)
  (dimension ApplicationStatus ?statusCode)
  (dimension DatabaseStatus ?dbStatusCode)
  (AvergeResponseTime (average ?elapsedTime) )
  (AverageDbQueryTime (average ?dbElapsedTime) )
  (NinetyPercentResponseTime (90-percentile ?elapsedTime) )
  (MaxResponseTime (max ?elapsedTime) )
  (TotalBytesSent (sum ?bytesSent) )
)
\end{lstlisting}

Each source begins with \fn{<-} followed by an ordinary pattern. Variables
shared between sources join them, exactly as in a rule: here the two logs
are joined on \fn{?transactionId}, and predicate constraints keep only
database queries that ran within the request's time window. A source may
use any constraint a rule pattern may. Several sources joined this way
form a star or snowflake schema: a portfolio cube, for instance, can join
positions to their accounts, to the stocks they hold and to the stocks'
ratings, and offer dimensions taken from all four.

A dimension is \fn{(dimension name ?variable)}: the name is what rules
will query by, the variable says where the value comes from. A measure is
\fn{(name (function ?variable))}, where the function is one of the
registered measures and the variable must be bound to a numeric slot. The
built-in measures are \fn{sum}, \fn{average}, \fn{median}, \fn{min},
\fn{max}, \fn{stdev} and the \fn{70-}, \fn{80-} and \fn{90-percentile};
\fn{(measures)} lists them, and \fn{load-measure} or
\fn{load-measure-group} add Java classes implementing
\fn{org.morendo.rete.measures.Measure}. Load such a measure before
defining any cube that uses it: \fn{defcube} rejects an unknown measure
with the message \fn{Measure \dots\ does not exist} and answers
\fn{false}.

\section{What the engine does with it}

Compiling a cube does four things: it validates the definition against the
templates, it creates a \emph{cube template} whose slots are the
dimensions and measures, it asserts one \emph{cube fact} of that template,
and it generates and compiles a rule that feeds the cube. \fn{ppdefcube}
shows the compiled cube, and \fn{(rules)} shows the generated rule, named
\fn{generated\_cube\_rule\_}\emph{cube}. Its conditions are the sources;
its action is \fn{(cube-add-data "TransactionCube")}, and its
modification action (Section~\ref{sec:modification-logic}) is
\fn{(cube-delete-data "TransactionCube")}. The rule is declared
\fn{no-agenda}, so the cube is updated the moment a matching combination
of facts appears or disappears, before any ordinary rule fires. Nothing in
the user's rule set needs to sequence it.

A cube is therefore \emph{reactive}: every assertion, modification or
retraction of a source fact changes the cube's contents immediately, and
the cube fact is re-asserted so that rules depending on it are
re-evaluated.

\section{Querying a cube from a rule}

The \fn{cubequery} conditional element slices the cube. Its pattern names
the cube and gives dimensions to filter by and measures to read;
\file{samples/molap/cube\_rules.clp}:

\begin{lstlisting}
(defrule ContentStats
  ?report <- (Report (Application ?app))
  ?cresult <- (cubequery
    (TransactionCube
      (Application ?app)
      (AvergeResponseTime ?aveRespTime)
      (MaxResponseTime ?maxRespTime)
      (AverageDbQueryTime ?aveDbTime)
      (NinetyPercentResponseTime ?ninetyPercent)
    )
  )
=>
  (printout t "results: " ?app " - ave resp: " ?aveRespTime
    " - max resp: " ?maxRespTime " - 90% resp: " ?ninetyPercent crlf)
  (printout t "- db ave: " ?aveDbTime " - db max: " ?maxDbTime crlf)
)
\end{lstlisting}

A dimension bound to a variable that is already bound earlier in the rule
(\fn{?app}) selects the slice; a measure bound to a fresh variable
receives the aggregate over that slice. The rule fires once per
\fn{Report} with the measures for that application, and fires again when
the cube changes. The result of a query is cached, so several rules
asking for the same slice share one computation.

The compiler switches indexing on for every dimension a \fn{cubequery}
uses; other dimensions are not indexed. \fn{(index-dimension cube dim
...)} turns indexing on by hand, for dimensions queried only from Java.

\section{Inspecting and profiling}

\begin{center}
\begin{tabularx}{\textwidth}{>{\raggedright\arraybackslash}p{2.9in}L}
\toprule
\fn{(list-defcubes)}, \fn{(cubes)} & names of the cubes \\
\fn{(ppdefcube name)} & dimensions with their index state and size, and the measures \\
\fn{(profile-cube name)}, \fn{(unprofile-cube name)} & time the cube's queries; \fn{print-profile} reports the average \\
\fn{(profile-cube-index name)}, \fn{(print-profile-cube-index)} & statistics on the dimension indexes \\
\bottomrule
\end{tabularx}
\end{center}

The scenario \file{src/test/resources/scenarios/molap.clp} loads the
sample's templates, cube and rule; the repository ships no log data for
it, so to see the rule fire, assert matching \fn{HTTPLog},
\fn{DatabaseQueryLog} and \fn{Report} facts of your own. The templates in
\file{samples/molap/templates.clp} type their time slots as \fn{DATE},
which the time functions in the cube's predicates require.
